\reading{MR-2}{Non-Parametric Approaches}{strike}{4}
% ==========================================================

\begin{mrkeybox}[title={Learning objectives in this reading}]
\small
\begin{enumerate}[label=\textbf{\alph*.},leftmargin=1.7em]
\item Apply the bootstrap historical simulation approach to estimate coherent risk measures.
\item Describe historical simulation using non-parametric density estimation.
\item Compare and contrast the age-weighted, the volatility-weighted, the correlation-weighted,
      and the filtered historical simulation approaches.
\item Identify advantages and disadvantages of non-parametric estimation methods.
\end{enumerate}
\end{mrkeybox}

\section{Basics}
\begin{itemize}
\item \term{Parametric methods:} which assume specific distributions like normal or lognormal.
\item \term{Non-parametric estimation and bootstrapping} are data-driven approaches that do not
      rely on specific assumptions about the underlying distribution. These methods are
      particularly useful for VaR estimation when rare events occur in the tail of the
      distribution.
\end{itemize}

% ---------------------------------------------------------
\lo{MR-2 a}{Apply the bootstrap historical simulation approach to estimate coherent risk measures.}

\term{Bootstrap historical simulation.} This is a straightforward method for estimating VaR.

\begin{enumerate}[label=\textbf{\alph*)},leftmargin=1.9em]
\item \term{VaR estimation:} it involves repeatedly sampling from the original data set,
      calculating VaR for each sample, and returning the data. This process is done
      \textit{with replacement}, meaning the same data points can be selected multiple times.
      The best VaR estimate from the full data set is the average of all sample VaRs.
\item \term{ES estimation:} the same process can be applied to estimate expected shortfall (ES),
      where the tail region is divided into slices and VaRs are averaged at each slice. The best
      estimate of ES for the original data set is the average of all the sample expected
      shortfalls.
\end{enumerate}

\begin{center}
\begin{tikzpicture}[every node/.style={font=\scriptsize\sffamily}]
% ---- parent distribution ----
\begin{scope}[shift={(2.0,0)}]
  \draw[FRMNavy, line width=0.9pt, fill=PaleNavy]
    plot[smooth, domain=0:6, samples=100]
    (\x, {1.35*exp(-((\x-2.1)^2)/1.7) + 0.98*exp(-((\x-4.3)^2)/1.3)})
    -- (6,0) -- (0,0) -- cycle;
  \draw[black!60, line width=0.6pt] (-0.1,0) -- (6.1,0);
  \node[anchor=west, text=black!65] at (6.35,0.75) {Original data set};
\end{scope}
% ---- four resamples ----
\foreach \i/\x/\v/\p/\q in {1/0.0/10/1.8/4.4, 2/2.6/5/2.4/4.0, 3/5.2/7/2.0/4.6, 4/7.8/8/2.5/4.2}{
  \begin{scope}[shift={(\x,-2.6)}]
    \draw[FRMNavy!85, line width=0.7pt, fill=PaleNavy!70]
      plot[smooth, domain=0:6, samples=80]
      ({\x*0.33}, 0);
    \draw[FRMNavy!85, line width=0.7pt, fill=PaleNavy!70]
      plot[smooth, domain=0:6, samples=80]
      ({\x*0.33}, {0.42*(1.30*exp(-((\x-\p)^2)/1.5) + 0.95*exp(-((\x-\q)^2)/1.1))})
      -- (1.98,0) -- (0,0) -- cycle;
    \draw[black!55, line width=0.5pt] (-0.05,0) -- (2.03,0);
  \end{scope}
  \node[text=black!80] at (\x+0.99,-3.0) {VaR = \v\,M\$};
}
% ---- arrows ----
\foreach \x in {0.99,3.59,6.19,8.79}{
  \draw[FRMNavy!45, -{Stealth[length=4.5pt]}, line width=0.6pt]
    (4.9,-0.18) .. controls (4.9,-1.1) and (\x,-1.3) .. (\x,-2.42);
}
\node[draw=FRMNavy, fill=PaleNavy, rounded corners=2pt, inner sep=5pt, text=black]
  at (4.9,-3.85) {\textbf{Bootstrapping VaR} $=\;\dfrac{10+5+7+8}{4}\;=\;7.5$ M\$};
\end{tikzpicture}
\end{center}
\figcap{One original data set is resampled four times with replacement. Each resample produces
its own VaR, and the bootstrap estimate is the simple average of those four figures.}

\begin{mrkeybox}
Empirical analysis shows that bootstrapping provides more precise estimates of coherent risk
measures compared to historical simulation alone.
\end{mrkeybox}

% ---------------------------------------------------------
\lo{MR-2 b}{Describe historical simulation using non-parametric density estimation.}

The clear advantage of the traditional historical simulation approach is its \term{simplicity}.
One obvious \term{drawback}, however, is that the discreteness of the data does not allow for
estimation of VaRs between data points (e.g.\ 95.5\% VaR).

\begin{mrkeybox}[title={Advantages of non-parametric density estimation}]
\begin{itemize}
\item Does not rely on restrictive assumptions about the underlying distribution.
\item Allows for VaR calculation at various confidence levels.
\end{itemize}
\end{mrkeybox}

\textbf{VaR calculation for 95.5\%.} Connecting the midpoints between histogram bars in the data
set creates a \term{surrogate density function}. This non-parametric approach improves upon
traditional historical simulation by enabling VaR calculation for a continuum of points in the
data set.

\begin{center}
\begin{tikzpicture}[every node/.style={font=\scriptsize\sffamily}]
% (a) original histogram
\begin{scope}
  \fill[FRMSoft, draw=black!55] (0,0) rectangle (1.1,1.0);
  \fill[FRMSoft, draw=black!55] (1.1,0) rectangle (2.2,1.8);
  \fill[FRMSoft, draw=black!55] (2.2,0) rectangle (3.3,2.5);
  \draw[FRMNavy, line width=1.1pt] (0.55,1.0) -- (1.65,1.8) -- (2.75,2.5) -- (3.3,2.62);
  \node[anchor=south, text=black!70] at (1.65,3.0) {(a) Original histogram};
\end{scope}
% (b) surrogate density
\begin{scope}[shift={(5.0,0)}]
  \fill[FRMSoft, draw=black!45, dotted] (0.55,0) rectangle (1.65,1.0);
  \fill[FRMSoft, draw=black!45, dotted] (1.65,0) rectangle (2.75,1.8);
  \fill[FRMSoft, draw=black!45, dotted] (2.75,0) rectangle (3.85,2.5);
  \fill[black!25, draw=black!60] (0,0) rectangle (0.55,1.0);
  \fill[FRMNavy!30, draw=none] (0.55,1.0) -- (1.10,1.40) -- (1.10,1.0) -- cycle;
  \draw[FRMNavy, line width=1.1pt] (0.55,1.0) -- (1.65,1.8) -- (2.75,2.5) -- (3.3,2.62);
  \draw[black!60] (0,0) -- (3.85,0);
  \node[anchor=south, text=black!70] at (1.9,3.0) {(b) Surrogate density function};
\end{scope}
\end{tikzpicture}
\end{center}
\figcap{Panel (a) is the raw histogram, where VaR can only be read off at a bar edge. Panel (b)
joins the bar midpoints into a continuous surrogate density, so a confidence level falling
between two bars --- 95.5\%, for instance --- now has a value.}

% ---------------------------------------------------------
\lo{MR-2 c}{Compare and contrast the age-weighted, the volatility-weighted, the correlation-weighted, and the filtered historical simulation approaches.}

The previous weighted historical simulation, discussed in \term{MR-1}, assumed that both current
and past (arbitrary) $n$ observations up to a specified cutoff point are used when computing the
current period VaR.

\begin{itemize}
\item Older observations beyond the cutoff date are assumed to have a zero weight and the
      relevant $n$ observations have equal weight of $1/n$.
\item \term{``Ghost effects''} of past events can persist in the computation of current VaR
      until they disappear after $n$ periods.
\end{itemize}

This reading identifies \term{four improvements} to the traditional historical simulation method.

\subsection{1) Age-weighted historical simulation}
One adjustment to the equal-weighted assumption used in historical simulation is to weight
recent observations more and distant observations less.

Assume $w(1)$ is the probability weight for the observation that is one day old. Then $w(2)$ can
be defined as $\lambda w(1)$, $w(3)$ can be defined as $\lambda^{2} w(1)$, and so on.

\begin{center}
\begin{tikzpicture}[every node/.style={font=\scriptsize\sffamily}]
  \node[anchor=west, text=black] at (0,0.66) {\textbf{Dec 2021}};
  \node[anchor=east, text=black] at (15.6,0.66) {\textbf{Dec 2022}};
  \draw[FRMNavy, line width=0.9pt] (0,0.32) -- (15.6,0.32);
  \foreach \x/\l in {0.35/0, 0.95/1, 1.55/2, 2.15/3}{
    \node[text=FRMGold] at (\x,0.00) {\l};}
  \node[text=FRMGold] at (2.95,0.00) {$\cdots$};
  \node[text=black!40] at (6.4,0.00) {$\cdots\qquad\cdots\qquad\cdots\qquad\cdots$};
  \node[text=FRMGold] at (10.1,0.00) {$\cdots$};
  \foreach \x/\l in {11.60/249, 12.95/250, 14.30/251, 15.35/252}{
    \node[text=FRMGold] at (\x,0.00) {\l};}
  \node[text=black] at (11.60,-0.72) {$\lambda^{3}\omega_{(1)}$};
  \node[text=black] at (12.95,-0.72) {$\lambda^{2}\omega_{(1)}$};
  \node[text=black] at (14.30,-0.72) {$\lambda\,\omega_{(1)}$};
  \node[text=black] at (15.35,-0.72) {$\omega_{(1)}$};
\end{tikzpicture}
\end{center}
\figcap{The 252 observations in the window, oldest on the left and newest on the right. The most
recent observation carries the full weight $\omega_{(1)}$ and each step back in time multiplies
that weight by another $\lambda$.}

\begin{mrfmlbox}
\[ w(i) \;=\; \frac{\lambda^{\,i-1}\,(1-\lambda)}{1-\lambda^{\,n}} \]
\mrvk{$w(i)$ = probability weight given to the observation that is $i$ days old;\quad
$\lambda$ = decay parameter, $0 < \lambda \le 1$;\quad $n$ = number of observations in the
sample.}{}
\end{mrfmlbox}

\begin{center}
\begin{tikzpicture}
\begin{axis}[
  width=11.0cm, height=5.2cm,
  xmin=0, xmax=255, ymin=0, ymax=5.6,
  xlabel={\scriptsize Age of the observation, $i$ (days)},
  ylabel={\scriptsize Weight $w(i)$ (\%)},
  xlabel style={font=\scriptsize}, ylabel style={font=\scriptsize},
  tick label style={font=\scriptsize},
  xtick={0,50,100,150,200,250}, ytick={0,1,2,3,4,5},
  legend style={font=\scriptsize\sffamily, draw=FRMRule, fill=white,
    at={(0.5,-0.30)}, anchor=north, legend columns=-1, column sep=8pt},
  axis line style={draw=black!55}, clip=false]
\addplot[FRMRed, line width=1.3pt, smooth] coordinates {
(1,5.000)(5,4.073)(10,3.151)(20,1.887)(30,1.130)(45,0.523)(60,0.243)(80,0.087)
(100,0.031)(130,0.007)(160,0.001)(200,0.000)(252,0.000)};
\addlegendentry{$\lambda = 0.95$ --- fast decay}
\addplot[FRMNavy, line width=1.3pt, dashed, smooth] coordinates {
(1,1.086)(5,1.044)(10,0.992)(20,0.898)(30,0.812)(45,0.698)(60,0.600)(80,0.491)
(100,0.402)(130,0.297)(160,0.220)(200,0.147)(252,0.087)};
\addlegendentry{$\lambda = 0.99$ --- slow decay}
\addplot[FRMGreen, line width=1.3pt, dotted] coordinates {(1,0.397)(252,0.397)};
\addlegendentry{$\lambda = 1$ --- no decay, plain HS}
\node[anchor=west, font=\tiny\sffamily, text=FRMGreen, fill=white, inner sep=1.2pt]
  at (axis cs:105,1.45) {$\lambda=1$: every observation weighted $1/n = 0.397\%$};
\draw[FRMGreen!70, -{Stealth[length=3pt]}, line width=0.5pt]
  (axis cs:150,1.35) -- (axis cs:155,0.47);
\end{axis}
\end{tikzpicture}
\end{center}
\figcap{The same 252 observations under three decay parameters. At $\lambda = 0.95$ the weight
has effectively vanished by day 60, so the VaR is being computed from about two months of data
no matter how long the window is. At $\lambda = 0.99$ the oldest observation still carries 8\%
of the newest one's weight. At $\lambda = 1$ the line is flat and the method collapses back to
ordinary historical simulation. This is what ``close to 1 means slow decay'' actually looks
like.}

\begin{mrkeybox}[title={Remember --- the decay parameter}]
The $\lambda$ term (decay parameter) is between 0 and 1, and reflects the exponential rate of
decay in the weight or value given to an observation as it ages:
\begin{itemize}
\item a $\lambda$ \term{close to 1} indicates a slow rate of decay;
\item a $\lambda$ \term{far away from 1} indicates a high rate of decay, so only the most very
      recent data set is used;
\item a $\lambda$ \term{of 1} means no decay, so all historic data is used --- just like
      historical data simulation.
\end{itemize}
\end{mrkeybox}

\subsection{2) Volatility-weighted historical simulation}
The volatility-weighted historical simulation approach weights individual observations based on
their volatility instead of their proximity to the current date.

The intuition is that if recent volatility has \term{increased}, then using historical data will
\term{underestimate} the current risk level. Similarly, if current volatility is markedly
\term{reduced}, the impact of older data with higher periods of volatility will \term{overstate}
the current risk level.

\begin{mrfmlbox}
\[ r^{*}_{t,i} \;=\; \left( \frac{\sigma_{T,i}}{\sigma_{t,i}} \right) r_{t,i}
   \qquad\longrightarrow\qquad
   \frac{\text{current volatility estimate}}{\text{volatility estimate at } t} \times r_t \]
\mrvk{$r_{t,i}$ = the actual historical return on asset $i$ on day $t$;\quad
$r^{*}_{t,i}$ = the volatility-adjusted return;\quad
$\sigma_{T,i}$ = the current (most recent) volatility forecast for asset $i$;\quad
$\sigma_{t,i}$ = the volatility forecast for asset $i$ that applied on day $t$.}{}
\end{mrfmlbox}

Thus, the volatility-adjusted return is replaced with a larger (smaller) expression if current
volatility exceeds (is below) historical volatility on day $i$. Now, VaR, ES, and any other
coherent risk measure can be calculated in the usual way after substituting historical returns
with volatility-adjusted returns.

\begin{mrkeybox}[title={Advantages of the volatility-weighted method}]
\begin{enumerate}
\item \term{Explicitly incorporates volatility:} the method takes into account the volatility of
      the market, which provides a more accurate estimation compared to other historical methods
      that do not explicitly consider volatility.
\item \term{Sensible near-term estimates:} the method produces VaR estimates that are more
      sensible and relevant in relation to current market conditions, allowing for better risk
      assessment.
\item \term{Higher VaR estimates:} the volatility-adjusted returns enable VaR estimates that are
      higher than those obtained from using the historical data set alone, offering a more
      conservative measure of risk.
\end{enumerate}
\end{mrkeybox}

\subsection{3) Correlation-weighted historical simulation}
\begin{itemize}
\item Correlation-weighting is a little more involved than volatility-weighting.
\item To see the principles involved, suppose for the sake of argument that we have already made
      any volatility-based adjustments to our HS returns along Hull-White lines, but also wish to
      adjust those returns to reflect changes in correlations.
\end{itemize}

\subsection{4) Filtered historical simulation (FHS)}
Combines the historical simulation model with a GARCH or AGARCH model.

\textbf{The steps for FHS are as follows:}
\begin{enumerate}
\item Firstly, use the historical return to find any surprise and thus reproduce volatility with
      a GARCH or AGARCH model.
\item Secondly, these volatility forecasts are then divided into the realized returns to produce
      a set of standardized returns, which is i.i.d.
\item The third stage involves bootstrapping from the set of standardized returns.
\item Finally, each of these simulated returns gives us a possible end-of-tomorrow portfolio
      value, and a corresponding possible loss, and we take the VaR to be the loss corresponding
      to our chosen confidence level.
\end{enumerate}


\subsection{Summary: the four approaches compared}

\begin{center}
\small
\begin{tabularx}{\linewidth}{@{}p{2.35cm} >{\raggedright\arraybackslash}X >{\raggedright\arraybackslash}X@{}}
\arrayrulecolor{FRMRule}\toprule
\rowcolor{FRMNavy} \hdr{Approach} & \hdr{What it re-weights, and how} & \hdr{Key point}\\
\midrule
\textbf{Age-weighted}
  & Weights observations by \textit{how old they are}. The newest observation gets
    $\omega_{(1)}$ and each further day back is multiplied by $\lambda$.
  & $\lambda$ close to 1 decays slowly; $\lambda$ far from 1 decays fast; $\lambda = 1$
    collapses back to plain historical simulation.\\
\rowcolor{FRMSoft}
\textbf{Volatility-\newline weighted}
  & Weights observations by \textit{their volatility}, not their proximity to today. Each past
    return is scaled by $\sigma_{T,i} / \sigma_{t,i}$.
  & Raises past returns when current volatility is higher than it was then, and lowers them when
    it is lower. Produces higher, more conservative VaR estimates.\\
\textbf{Correlation-\newline weighted}
  & Adjusts returns to reflect \textit{changes in correlations}, on top of any volatility
    adjustment already made along Hull-White lines.
  & More involved than volatility-weighting.\\
\rowcolor{FRMSoft}
\textbf{Filtered \newline (FHS)}
  & Combines historical simulation with a \textit{GARCH or AGARCH} model: standardise the
    returns by the volatility forecast, then bootstrap from the standardised set.
  & The standardised returns are i.i.d., which is what makes the bootstrap step valid.\\
\bottomrule
\end{tabularx}
\end{center}
\figcap{The four differ only in what they re-weight: age-weighting uses time, volatility-weighting
uses the volatility ratio, correlation-weighting uses the correlation structure, and filtered
historical simulation replaces the raw returns with model-standardised ones before resampling.}

% ---------------------------------------------------------
\lo{MR-2 d}{Identify advantages and disadvantages of non-parametric estimation methods.}

\begin{center}
\small
\begin{tabularx}{\linewidth}{@{}>{\raggedright\arraybackslash}X >{\raggedright\arraybackslash}X@{}}
\arrayrulecolor{FRMRule}\toprule
\rowcolor{FRMNavy} \hdr{Advantages} & \hdr{Disadvantages}\\
\midrule
a.~Intuitive and conceptually simple.
  & a.~Very dependent on the historical data set.\\
\rowcolor{FRMSoft}
b.~Do not depend on parametric assumptions.
  & b.~Subject to ghost effect.\\
c.~Accommodate any type of position.
  & c.~If our data period was unusually quiet, non-parametric methods will often produce VaR or
      ES estimates that are too low, and vice versa.\\
\rowcolor{FRMSoft}
d.~No need for covariance matrices, no curses of dimensionality.
  & d.~Have difficulty (act slowly) handling shifts (permanent risk change) that take place
      during our sample period.\\
e.~Use data that are (often) readily available.
  & \\
\bottomrule
\end{tabularx}
\end{center}


% ==========================================================
